\chapter{Laparoscopic subtotal hysterectomy}
\index{Laparoscopic subtotal hysterectomy}

\section*{Definition and Overview}
Laparoscopic subtotal hysterectomy (LSH) is a minimally invasive surgical procedure in which the uterine corpus (body of the uterus) is removed while the cervix is preserved, performed using laparoscopic (keyhole) techniques. LSH is commonly indicated for benign gynaecological conditions such as symptomatic uterine fibroids (leiomyomas), adenomyosis, pelvic endometriosis, and abnormal uterine bleeding that has not responded to medical therapy. By retaining the cervix, the procedure aims to preserve pelvic floor integrity, support pelvic organs, and maintain sexual function, although comparative clinical trial data regarding pelvic floor outcomes between subtotal and total hysterectomies remains mixed. This minimally invasive approach offers significant clinical benefits over open abdominal hysterectomy, including reduced postoperative pain, lower blood loss, shorter hospital stay, and a faster return to baseline activities.

\section*{Epidemiology}
Hysterectomy is one of the most frequently performed major gynaecological procedures worldwide. In many countries and other developed countries, thousands of hysterectomies are performed annually, with a steady clinical transition from open laparotomy to minimally invasive routes, including laparoscopic, vaginal, and robotic approaches. Subtotal hysterectomy accounts for approximately 10\% to 20\% of all hysterectomies performed for benign indications. The peak age group for undergoing LSH is between 40 and 50 years, correlating with the highest prevalence of symptomatic uterine fibroids and adenomyosis. The choice of laparoscopic subtotal hysterectomy depends on the patient's individual clinical profile, uterine size, and absence of cervical pathology, as it is strictly contraindicated in patients with a history of cervical dysplasia or suspected gynaecological malignancy.

\section*{Aetiology and Pathophysiology}
Laparoscopic subtotal hysterectomy is indicated for benign conditions of the uterus that cause chronic pain, pressure symptoms, or severe haemorrhage. The physiological basis of the procedure involves the laparoscopic dissection of the uterine ligaments and vascular supply, followed by transection of the uterine body at the level of the internal cervical os. The underlying benign pathologies that lead to LSH include:
\begin{itemize}
    \item \textbf{Uterine fibroids (leiomyomas):} Benign clonal neoplasms of smooth muscle cells within the myometrium, which cause uterine enlargement, pelvic pressure, dysmenorrhoea, and menorrhagia due to increased endometrial vascularity and surface area.
    \item \textbf{Adenomyosis:} The abnormal infiltration of endometrial glands and stroma into the myometrium, resulting in reactive myometrial hypertrophy, global uterine enlargement, severe dysmenorrhoea, and menorrhagia.
    \item \textbf{Endometriosis:} The presence of ectopic endometrial tissue outside the uterine cavity, leading to chronic pelvic inflammation, anatomical distortions, pelvic adhesions, and cyclic or acyclic pelvic pain.
    \item \textbf{Refractory abnormal uterine bleeding:} Menstrual bleeding of excessive duration or volume that has failed to improve with pharmacological treatments (such as hormonal therapies or antifibrinolytics) or less invasive procedures (such as endometrial ablation).
\end{itemize}

\section*{Clinical Features}
The clinical features of patients undergoing LSH are primarily related to the preoperative benign pathologies, followed by postoperative signs of healing. Clinical features include:
\begin{itemize}
    \item \textbf{Menorrhagia:} Excessive, prolonged menstrual bleeding, often with clots, leading to severe fatigue and iron-deficiency anaemia.
    \item \textbf{Pelvic mass and bulk symptoms:} Progressive abdominal distension, pelvic heaviness, urinary frequency from bladder compression, or constipation from rectal pressure.
    \item \textbf{Chronic pelvic pain and dysmenorrhoea:} Severe pain during menstruation or persistent lower abdominal discomfort that impairs daily functioning.
    \item \textbf{Dyspareunia:} Pain during sexual intercourse, often associated with adenomyosis or co-existing pelvic endometriosis.
    \item \textbf{Postoperative pain:} Mild-to-moderate incisional pain at laparoscopic port sites and referred shoulder tip pain, which is caused by residual carbon dioxide gas irritating the diaphragm.
    \item \textbf{Cyclic spotting:} Continued light, periodic vaginal bleeding in approximately 5\% to 10\% of patients who retain a small amount of active endometrium at the cervical stump.
\end{itemize}

\section*{Differential Diagnosis}
Before proceeding with LSH, other potential sources of abnormal uterine bleeding, pelvic pain, or a pelvic mass must be excluded. The differential diagnosis includes:
\begin{itemize}
    \item \textbf{Gynaecological malignancy:} Endometrial carcinoma, cervical cancer, or uterine sarcoma, which must be ruled out as they require total hysterectomy and oncological staging, and contraindicate LSH due to the risk of intraperitoneal dissemination of malignant cells during tissue morcellation.
    \item \textbf{Pelvic inflammatory disease (PID):} Chronic upper genital tract infection causing pelvic pain and adhesions, which must be differentiated from endometriosis or adenomyosis.
    \item \textbf{Ovarian lesions:} Benign ovarian cysts, endometriomas, or ovarian neoplasms that present with pelvic pressure, masses, or cyclic pain.
    \item \textbf{Systemic bleeding disorders:} Inherited coagulopathies, such as von Willebrand disease, presenting with menorrhagia that may be manageable without hysterectomy.
    \item \textbf{Endocrine disorders:} Thyroid dysfunction or hyperprolactinaemia causing hormonal imbalances and abnormal bleeding patterns.
\end{itemize}

\section*{When to Seek Medical Care}
Patients recovering from LSH should contact their gynaecologist or surgical team if they experience worsening pain, persistent low-grade fever, or unusual vaginal discharge.

\textbf{Contact your local emergency services for:}
\begin{itemize}
    \item \textbf{Haemorrhage:} Heavy, continuous vaginal bleeding or sudden soaking of multiple sanitary pads within an hour.
    \item \textbf{Severe pelvic or abdominal pain:} Sudden-onset, escalating pain that does not respond to prescribed analgesics, which may indicate internal bleeding or visceral injury.
    \item \textbf{Sepsis:} High fever (above 38.5 degrees Celsius), severe chills, shaking, rapid heart rate, or confusion.
    \item \textbf{Pulmonary embolism:} Sudden shortness of breath, sharp chest pain, or coughing up blood.
    \item \textbf{Deep vein thrombosis:} Unilateral calf or leg pain, swelling, warmth, or redness.
    \item \textbf{Bowel or bladder dysfunction:} Inability to pass urine, severe abdominal bloating, or persistent vomiting.
\end{itemize}

\section*{Diagnosis and Investigations}
A comprehensive diagnostic evaluation is essential to confirm benign disease, verify cervical health, and assess anaesthetic fitness before LSH. Key investigations include:
\begin{itemize}
    \item \textbf{Cervical screening test (CST):} A recent negative result for high-risk HPV and cytology is mandatory to ensure the safety of retaining the cervix.
    \item \textbf{Transvaginal ultrasound (TVUS):} The standard imaging modality to evaluate uterine dimensions, map the location and size of fibroids, and assess endometrial thickness and the ovaries.
    \item \textbf{Endometrial biopsy:} Indicated in patients with abnormal uterine bleeding, particularly those aged over 45 years, to rule out endometrial hyperplasia or atypical changes.
    \item \textbf{Full blood count (FBC):} To evaluate for iron-deficiency anaemia secondary to menorrhagia, guiding preoperative iron supplementation or transfusion.
    \item \textbf{Magnetic resonance imaging (MRI):} Utilised in complex cases to differentiate between fibroids and adenomyosis, or to map large masses prior to surgery.
\end{itemize}

\section*{Management}
The management of LSH involves preoperative counselling, precise intraoperative technique, and structured postoperative recovery protocols. Management components include:
\begin{itemize}
    \item \textbf{Surgical procedure:} Under general anaesthesia, laparoscopic access is obtained. The uterine body is mobilised by dividing the round ligaments, fallopian tubes, and ovarian ligaments (with preservation of the ovaries). The uterine corpus is transected from the cervix at the level of the internal os. The specimen is removed, often using power morcellation within a containment bag.
    \item \textbf{Pain management:} A multimodal analgesia regimen including paracetamol, non-steroidal anti-inflammatory drugs (NSAIDs), and oral opioids for breakthrough pain, combined with local anaesthetic infiltration at incision sites.
    \item \textbf{Thromboprophylaxis:} Early postoperative mobilisation, mechanical compression devices, and pharmacological agents (such as low-molecular-weight heparin) based on individual VTE risk assessments.
    \item \textbf{Ongoing cervical screening:} Counselling patients that because the cervix remains intact, they must continue routine cervical screening according to national guidelines.
    \item \textbf{Hormone replacement therapy (HRT):} Not routinely required postoperatively if the ovaries are preserved; however, if oophorectomy is also performed, HRT may be discussed for symptom control in premenopausal patients.
\end{itemize}

\section*{Complications}
While laparoscopic subtotal hysterectomy is generally safe, it carries risks of intraoperative and postoperative complications. These include:
\begin{itemize}
    \item \textbf{Vascular injury and haemorrhage:} Major bleeding from uterine or pelvic vessels that may necessitate blood transfusion or conversion to laparotomy.
    \item \textbf{Visceral injury:} Inadvertent thermal or mechanical damage to the bladder, ureters, or bowel, requiring surgical repair.
    \item \textbf{Infection:} Pelvic infection, vaginal vault haematoma infection, urinary tract infection, or surgical site infection at the laparoscopic port sites.
    \item \textbf{Venous thromboembolism:} Deep vein thrombosis or pulmonary embolism arising from pelvic surgery and immobility.
    \item \textbf{Cyclic bleeding:} Persistent light menstruation resulting from residual endometrial cells left at the cervix, which occurs in 5\% to 10\% of patients.
    \item \textbf{Cervical stump dysplasia:} The long-term risk of developing cervical intraepithelial neoplasia (CIN) or cervical cancer, necessitating continued surveillance.
\end{itemize}

\section*{Prognosis and Outcomes}
The long-term prognosis after LSH is excellent. The majority of patients report significant relief from preoperative symptoms such as menorrhagia and pelvic pressure, leading to a marked improvement in physical and psychological well-being. Because LSH is minimally invasive, the recovery phase is typically brief. Most patients are discharged within 24 to 48 hours and return to light duties within 2 to 3 weeks, and full activity within 4 to 6 weeks. The retention of the cervix is associated with high levels of patient satisfaction, although patient education regarding the potential for minor cyclic bleeding and the critical need for continued cervical screening is vital for optimal long-term outcomes.

\section*{Prevention and Self-Care}
Postoperative self-care is focused on preventing complications, supporting tissue healing, and ensuring a safe return to normal activities. Recommendations include:
\begin{itemize}
    \item \textbf{Graduated activity:} Engaging in short, regular walks starting on the day of surgery to prevent deep vein thrombosis and aid bowel function.
    \item \textbf{Avoiding abdominal strain:} Restricting heavy lifting (greater than 5 kilograms) and strenuous core exercises for at least 4 to 6 weeks to protect the pelvic floor and fascial incisions.
    \item \textbf{Pelvic rest:} Avoiding vaginal intercourse, tampons, and vaginal douching for 6 weeks to minimise the risk of pelvic infection and allow the cervical stump to heal.
    \item \textbf{Wound care:} Keeping laparoscopic port incisions clean and dry, and monitoring closely for redness, swelling, or drainage.
    \item \textbf{Dietary adjustments:} Consuming a high-fibre diet and drinking plenty of fluids to avoid constipation, which can strain pelvic tissues.
\end{itemize}

\section*{Sources and Further Reading}
The following reputable organisations provide further information on laparoscopic subtotal hysterectomy and gynaecological surgery:
\begin{itemize}
    \item Royal Australian and New Zealand College of Obstetricians and Gynaecologists (RANZCOG). \textit{Patient resource on hysterectomy options and recovery.} https://ranzcog.edu.au/
    \item Healthdirect Australia. \textit{Overview of laparoscopic hysterectomy procedures and post-surgical care.} https://www.healthdirect.gov.au/
    \item American College of Obstetricians and Gynecologists (ACOG). \textit{Clinical information and patient guides for hysterectomy.} https://www.acog.org/
    \item NHS. \textit{Information on subtotal and total hysterectomy surgery and recovery guidelines.} https://www.nhs.uk/
    \item Better Health Channel. \textit{Gynaecological surgical procedures and post-operative support services.} https://www.betterhealth.vic.gov.au/
\end{itemize}
